\section{Residual attachments by matching restriction}
\label{sec:residual-attachments-v3}

Section~8 bounded the sum of the bare high-skeleton weights. It did not include
the rewards of unexposed cells or the binary cycle-space factor of the residual
support. We now condition on one feasible high skeleton and estimate precisely
those remaining factors.

\subsection{Conditional decomposition}

Fix a canonical high skeleton with exposed block matching $M$, exposed
multiplicities $j$, and total exposed mass $J$. Put
\[
  m_0:=n-J.
\]
Let $(d_a)_a$ and $(d'_b)_b$ be the residual row and column degrees. Their two
sums equal $m_0$, and every degree is at most the phase cap $U$.

The unexposed pairs form a uniform bipartite configuration matching with these
residual degrees. Let $r'_{ab}$ be its cell counts and let
$H_{\mathrm{res}}$ be the simple support graph of cells with $r'_{ab}\ge2$.
The residual event $\mathcal E(M,j)$ imposes the cap
$r'_{ab}\le\lfloor U/2\rfloor$ and forbids any further pair in a cell of $M$.
Define
\begin{equation}
  \mathcal A(M,j)
  :=
  \mathbb E_{\mathrm{res}}\!\left[
    \left(\prod_{a,b}g(r'_{ab})\right)
    2^{\beta(M\cup H_{\mathrm{res}})}
    \mathbf 1_{\mathcal E(M,j)}
  \right].
  \label{eq:residual-attachment-v3}
\end{equation}
Only the residual configuration matching is random in this expectation.

The exact overlap decomposition is
\begin{equation}
  \frac{\mathbb E Z^2}{(\mathbb E Z)^2}
  =
  \sum_{(M,j)}w_{\mathrm{hi}}(M,j)\,\mathcal A(M,j),
  \label{eq:exact-attachment-decomposition-v3}
\end{equation}
where the sum ranges over the finite family of feasible canonical high
skeletons. In view of Proposition~\ref{prop:bare-high-skeleton-v3}, it remains
to bound $\mathcal A(M,j)$ uniformly in $(M,j)$.

\subsection{Threshold expansion and cell activities}

For a residual cell $(a,b)$ outside $M$, put
\begin{equation}
  \theta_{ab}:=
  \frac{\mathrm e\,d_ad'_b}{m_0}.
  \label{eq:theta-v3}
\end{equation}
Let
\[
  \Delta_x:=g(x)-g(x-1),
  \qquad
  R:=\left\lfloor\frac U2\right\rfloor,
\]
and define
\begin{equation}
  \lambda_{ab}
  :=
  \sum_{x=3}^{R}\Delta_x\frac{\theta_{ab}^x}{x!},
  \qquad
  q_{ab}
  :=
  \frac{\theta_{ab}^2}{2}+\lambda_{ab}.
  \label{eq:lambda-q-v3}
\end{equation}
On $M$, set $\lambda_{ab}=q_{ab}=0$.

We first record the configuration estimate used in the expansion. For a finite
array of nonnegative demands $x=(x_{ab})$, write
$x_a=\sum_bx_{ab}$, $x'_b=\sum_ax_{ab}$, and $X=\sum_{a,b}x_{ab}$. Markov's
inequality applied to the product of binomial cell counts gives
\[
  \mathbb P_{\mathrm{res}}(r'_{ab}\ge x_{ab}\text{ for all }a,b)
  \le
  \frac{
    \prod_a(d_a)_{x_a}\prod_b(d'_b)_{x'_b}
  }{
    (m_0)_X\prod_{a,b}x_{ab}!
  }.
\]
If $X\le m_0$, then $(m_0)_X\ge(m_0/\mathrm e)^X$; if $X>m_0$, the event is
empty. Hence
\begin{equation}
  \mathbb P_{\mathrm{res}}(r'_{ab}\ge x_{ab}\text{ for all }a,b)
  \le
  \prod_{a,b}\frac{\theta_{ab}^{x_{ab}}}{x_{ab}!}.
  \label{eq:prescribed-cell-residual-v3}
\end{equation}
This joint estimate is the reason that cells sharing a row or a column may be
expanded simultaneously.

For $0\le r\le R$,
\[
  g(r)=1+\sum_{x=3}^{R}\Delta_x\mathbf 1_{\{r\ge x\}}.
\]
Let $E_0$ be the set of all residual cell positions outside $M$, and let
$\mathcal C(M)$ be the family of even subsets of $M\cup E_0$. For
$F\in\mathcal C(M)$, define
\[
\begin{split}
  \Phi_F(r')
  :=
  &\prod_{e\in F\setminus M}
  \left(
    \mathbf 1_{\{r'_e\ge2\}}
    +
    \sum_{x=3}^{R}\Delta_x\mathbf 1_{\{r'_e\ge x\}}
  \right)\\
  &\times
  \prod_{e\in E_0\setminus F}
  \left(
    1+
    \sum_{x=3}^{R}\Delta_x\mathbf 1_{\{r'_e\ge x\}}
  \right).
\end{split}
\]
The first product contains the support threshold required when $e$ belongs to
the even set. On the event $\mathcal E(M,j)$, expansion of the cycle-space
cardinality and of all local rewards gives the exact identity
\[
  \left(\prod_{a,b}g(r'_{ab})\right)
  2^{\beta(M\cup H_{\mathrm{res}})}
  =
  \sum_{F\in\mathcal C(M)}\Phi_F(r').
\]

\begin{lemma}[Fixed even-set expansion]
\label{lem:fixed-even-set-expansion-v3}
For every $F\in\mathcal C(M)$,
\[
  \mathbb E_{\mathrm{res}}\!\left[
    \Phi_F(r')\mathbf 1_{\mathcal E(M,j)}
  \right]
  \le
  \prod_{e\in F\setminus M}q_e
  \prod_{e\in E_0\setminus F}(1+\lambda_e).
\]
\end{lemma}

\begin{proof}
Expand $\Phi_F$ into its nonnegative threshold monomials. In a cell of
$F\setminus M$, a monomial chooses either the threshold-two term or one higher
increment; it never chooses both. In a cell outside $F\cup M$, it chooses
nothing or one higher increment. Apply
\eqref{eq:prescribed-cell-residual-v3} once to the complete demand of each
monomial. The threshold-two contribution is $\theta_e^2/2$, the higher
increments sum to $\lambda_e$, and the empty choice contributes one. Dropping
the cap and no-return indicator only enlarges the nonnegative expectation.
\end{proof}

Summing over $F$ and then inserting the missing factors
$1+\lambda_e\ge1$ gives
\begin{equation}
  \mathcal A(M,j)
  \le
  \left(\prod_{e\in E_0}(1+\lambda_e)\right)
  \sum_{F\in\mathcal C(M)}
    \prod_{e\in F\setminus M}q_e.
  \label{eq:lambda-even-factorization-v3}
\end{equation}

\subsection{Restriction outside the exposed matching}

The next finite lemma is independent of random graphs.

\begin{lemma}[Restriction-product bound]
\label{lem:restriction-product-v3}
Let $E$ be a finite set, let $\mathcal C$ be a finite family of subsets of
$E$, and let $I\subseteq E$. Suppose that
\[
  A\longmapsto A\setminus I
\]
is injective on $\mathcal C$. For nonnegative activities $(q_e)_{e\in E}$,
\[
  \sum_{A\in\mathcal C}
    \prod_{e\in A\setminus I}q_e
  \le
  \prod_{e\in E\setminus I}(1+q_e).
\]
\end{lemma}

\begin{proof}
The restrictions $A\setminus I$ form a subfamily of the power set of
$E\setminus I$ and occur without repetition. Enlarge the sum to the full power
set and expand the finite product.
\end{proof}

Apply the lemma with $E=M\cup E_0$, $I=M$, and
$\mathcal C=\mathcal C(M)$. The required injectivity is immediate but
important. If two even sets have the same restriction outside $M$, their
symmetric difference is an even subset of the matching $M$. Every nonempty
subset of a matching has a vertex of degree one, so this symmetric difference
must be empty. Therefore
\begin{equation}
  \sum_{F\in\mathcal C(M)}
    \prod_{e\in F\setminus M}q_e
  \le
  \prod_{e\in E_0}(1+q_e).
  \label{eq:even-family-product-v3}
\end{equation}

Since $0\le\lambda_e\le q_e$, equations
\eqref{eq:lambda-even-factorization-v3} and
\eqref{eq:even-family-product-v3}, together with $1+x\le e^x$, imply
\begin{equation}
  \mathcal A(M,j)
  \le
  \exp\!\left(2\sum_{e\in E_0}q_e\right).
  \label{eq:q-only-attachment-v3}
\end{equation}
Thus both the local rewards and the cycle-space cardinality are controlled by
one total residual activity.

\begin{remark}
Lemma~\ref{lem:restriction-product-v3} applies to any weighted set family whose
restriction map is injective. Here the family happens to be a binary cycle
space and injectivity follows solely from the fact that the deleted edge set is
a matching.
\end{remark}

\subsection{The intrinsic residual regime}

Assume
\begin{equation}
  2^U\le m_0^3.
  \label{eq:intrinsic-residual-regime-v3}
\end{equation}
Then $m_0\ge2^{U/3}$ and, by the degree cap,
\[
  \theta_{ab}
  \le
  \mathrm e U^2 2^{-U/3}.
\]

\begin{lemma}[Quadratic activity bound]
\label{lem:q-quadratic-v3}
There is an absolute constant $C_0$ such that, under
\eqref{eq:intrinsic-residual-regime-v3},
\begin{equation}
  q_{ab}\le C_0\theta_{ab}^2
  \label{eq:q-quadratic-v3}
\end{equation}
for every residual cell.
\end{lemma}

\begin{proof}
Fix a cell and write $\theta=\theta_{ab}$. The assertion is trivial when
$\theta=0$. Since $0\le\Delta_x\le g(x)$, it is enough to bound
\[
  a_x:=g(x)\frac{\theta^x}{x!},
  \qquad 3\le x\le R.
\]
For $x\ge3$,
\[
  \frac{a_{x+1}}{a_x}
  =
  \frac{2^x\theta}{x+1},
  \qquad
  \frac{a_{x+2}/a_{x+1}}{a_{x+1}/a_x}
  =
  2\frac{x+1}{x+2}>1.
\]
Thus $(a_x)$ is log-convex, and its maximum on $[3,R]$ occurs at an endpoint.
Consequently
\[
  \lambda_{ab}
  \le
  R(a_3+a_R).
\]
Now $a_3=(2/3)\theta^3$. Since
$R\theta=O(U^3 2^{-U/3})$, one has $Ra_3\le\theta^2$ for all sufficiently
large $U$.

For the other endpoint, using
$R=\lfloor U/2\rfloor$ and
$\theta\le\mathrm e U^2 2^{-U/3}$ gives
\[
  \log_2\theta
  \le
  \log_2\mathrm e+2\log_2U-\frac U3.
\]
For sufficiently large $U$, one has $R\ge3$ and $g(R)=2^{\binom R2-1}$. Hence
\[
\begin{aligned}
  \log_2\!\left(\frac{a_R}{\theta^2}\right)
  &=
  \binom R2-1+(R-2)\log_2\theta-\log_2(R!)\\
  &\le
  \left(\frac{U^2}{8}+O(U)\right)
  -\left(\frac{U^2}{6}+O(U)\right)
  +O(U\log U)\\
  &=
  -\frac{U^2}{24}+O(U\log U).
\end{aligned}
\]
The two quadratic terms come respectively from the local reward $g(R)$ and
from the factor $\theta^{R-2}$; the factorial term only improves the upper
bound. Hence $Ra_R\le\theta^2$ for all sufficiently large $U$. It follows that
$\lambda_{ab}\le2\theta^2$ in that range. For the finitely many remaining
values of $U$, the quotient $q_{ab}/\theta^2$ is bounded on the compact interval
$0\le\theta\le\mathrm e U^2 2^{-U/3}$, with limit $1/2$ at $\theta=0$.
Enlarging the constant proves the lemma.
\end{proof}

The cell intensities satisfy the exact identity
\[
  \sum_{a,b}\theta_{ab}^2
  =
  \frac{\mathrm e^2}{m_0^2}
  \left(\sum_a d_a^2\right)
  \left(\sum_b(d'_b)^2\right).
\]
Since every degree is at most $U$ and both degree sums equal $m_0$,
\[
  \sum_a d_a^2\le Um_0,
  \qquad
  \sum_b(d'_b)^2\le Um_0.
\]
Lemma~\ref{lem:q-quadratic-v3} therefore gives
\begin{equation}
  \sum_{a,b}q_{ab}\le C_1U^2.
  \label{eq:total-q-v3}
\end{equation}
Together with \eqref{eq:q-only-attachment-v3}, this proves
\begin{equation}
  \mathcal A(M,j)\le\exp(CU^2)
  \qquad\text{if }2^U\le m_0^3.
  \label{eq:intrinsic-attachment-v3}
\end{equation}

\subsection{The complementary residual regime}

Suppose instead that $2^U>m_0^3$. Then
\begin{equation}
  m_0<2^{\lceil U/3\rceil}.
  \label{eq:small-residual-mass-v3}
\end{equation}
We may now discard every residual restriction. Since each residual cell count
is at most $U$,
\[
  \prod_{a,b}g(r'_{ab})
  \le
  2^{\sum_{a,b}\binom{r'_{ab}}2}
  \le
  2^{(U-1)m_0/2}.
\]
Moreover, every edge of $H_{\mathrm{res}}$ uses at least two residual pairs,
so $|E(H_{\mathrm{res}})|\le m_0/2$. Equivalently, the restriction argument
above gives
\[
  2^{\beta(M\cup H_{\mathrm{res}})}
  \le
  2^{|E(H_{\mathrm{res}})|}
  \le
  2^{m_0/2}.
\]
Multiplication yields
\begin{equation}
  \mathcal A(M,j)
  \le
  2^{Um_0/2}
  \qquad\text{if }2^U>m_0^3.
  \label{eq:complementary-attachment-v3}
\end{equation}

\begin{proposition}[Uniform residual attachment]
\label{prop:uniform-residual-attachment-v3}
There is an absolute constant $C$ such that every feasible canonical high
skeleton satisfies
\[
  \mathcal A(M,j)
  \le
  \begin{cases}
    \exp(CU^2),&2^U\le m_0^3,\\
    2^{Um_0/2},&2^U>m_0^3.
  \end{cases}
\]
\end{proposition}

The phase estimates give
\[
  U=O(\log n),
  \qquad
  2^U=\Theta\!\left(\frac{n^2}{(\log n)^2}\right).
\]
In the intrinsic regime, $CU^2=O((\log n)^2)$. In the complementary regime,
\eqref{eq:small-residual-mass-v3} gives
\[
  Um_0
  =
  O\!\left(n^{2/3}(\log n)^{1/3}\right).
\]
Both exponents are $o(n/(\log n)^4)$, uniformly in the high skeleton. Hence
there is a deterministic sequence $\eta_n\to0$ such that
\begin{equation}
  \mathcal A(M,j)
  \le
  \exp\!\left\{
    \eta_n\frac{n}{(\log n)^4}
  \right\}
  \label{eq:uniform-attachment-error-v3}
\end{equation}
for every feasible skeleton.

\begin{proposition}[Normalized signed second moment]
\label{prop:normalized-second-moment-v3}
There is a deterministic sequence $\Lambda_n$ satisfying
\[
  \Lambda_n=o\!\left(\frac{n}{(\log n)^4}\right)
\]
such that
\begin{equation}
  1
  \le
  \frac{\mathbb E Z^2}{(\mathbb E Z)^2}
  \le
  e^{\Lambda_n}.
  \label{eq:normalized-second-moment-v3}
\end{equation}
\end{proposition}

\begin{proof}
Insert \eqref{eq:uniform-attachment-error-v3} into the exact decomposition
\eqref{eq:exact-attachment-decomposition-v3} and factor out the uniform
attachment bound. Proposition~\ref{prop:bare-high-skeleton-v3} controls the
remaining high-skeleton sum. The lower bound follows from nonnegativity of the
variance.
\end{proof}

Paley--Zygmund now yields
\[
  \mathbb P(Z>0)
  \ge
  \frac{(\mathbb E Z)^2}{\mathbb E Z^2}
  \ge e^{-\Lambda_n}.
\]
Section~10 amplifies this rare seed to a high-probability cocoloring.
